
\section{借余切丛理解光滑簇}

本节中，我们将构造代数闭域上的簇的一些双有理不变量（例如\trGeometricGenus），为\trUnramified{}态射这一观念提供动机，证明簇「几乎处处光滑」，并初窥\trHodge{}理论。

\textbf{21.5.1. \trGeometricGenus{}与其他源自 $i$-形式的双有理不变量。}\quad 设 $X$ 是 $k$ 上的射影\trScheme。那么 $\mathrm{h}^0(X,\Omega_{X/k})$ 是 $X$ 的一个不变量。事实上，如下面的\mbox{练习~21.5.A} 所示，当 $X$ 光滑时，它与相关的一些不变量都是\emph{双有理不变量}。我们先定义\textbf{（相对）$i$-形式层}（sheaf of (relative) $i$-forms）%
$\Omega^i_{X/Y}:=\wedge^i\Omega_{X/Y}$。$\Omega^i_{X/Y}$ 的\trSection（在某个\trOpenSet{}上）称为\textbf{（相对）$i$-形式}（在该\trOpenSet{}上）。

\textbf{21.5.2.}\ \emph{笑话。}\quad 老麦当劳有一个形式；$e_i\wedge e_i=0$。%
\footnote{译注：此句戏仿英文童谣《老麦当劳有个农场》（\emph{Old MacDonald Had a Farm}）中的唱词「Old MacDonald had a farm, E--I--E--I--O」，原文「had a form; $e_i\wedge e_i=0$」与之谐音。}

\textbf{21.5.A.（练习：$\mathrm{h}^0(X,\Omega^i_{X/k})$ 是双有理不变量）}\quad 设 $X$ 与 $X'$ 是彼此双有理等价的不可约光滑射影 $k$-\trVariety。证明（对每个 $i$）：$\mathrm{H}^0(X,\Omega^i_{X/k})\cong\mathrm{H}^0(X',\Omega^i_{X'/k})$。（提示：取定双有理映射 $\phi\colon X\dashrightarrow X'$。由\mbox{练习~15.3.B}，$\phi$ 的\trDomainOfDefinition{} $U$ 的\trComplement{}余维数至少为 $2$。把诸 $i$-形式从 $X'$ 拉回到 $U$ 上，我们便得到一个映射 $\phi^*\colon\mathrm{H}^0(X',\Omega^i_{X'/k})\to\mathrm{H}^0(U,\Omega^i_{X/k})$。利用\mbox{\trAlgebraic\trHartogs{}引理~13.5.19} 以及 $\Omega^i$ 是\trLocallyFreeSheaf{}这一事实，证明该映射可以延拓为映射 $\phi^*\colon\mathrm{H}^0(X',\Omega^i_{X'/k})\to\mathrm{H}^0(X,\Omega^i_{X/k})$。若 $\psi\colon X'\dashrightarrow X$ 是逆有理映射，我们同样得到一个映射 $\psi^*\colon\mathrm{H}^0(X,\Omega^i_{X/k})\to\mathrm{H}^0(X',\Omega^i_{X'/k})$。证明 $\phi^*$ 与 $\psi^*$ 互为逆映射，办法是证明每个复合映射都在 $X$ 或 $X'$ 的某个稠密\trOpenSubset{}上为恒等。）

\textbf{21.5.3.}\ \emph{典范丛 $\mathscr{K}_X$ 与\trGeometricGenus{} $\mathrm{p}_g(X)$。}\quad 若 $X$ 是 $n$ 维光滑 $k$-簇，则可逆层（即\trLineBundle）$\det\Omega_{X/k}=\Omega^n_{X/k}$（「\trAlgebraicVolumeForm」层）具有特殊的重要性，它被称为\textbf{典范（可逆）层}，或者\textbf{典范（线）丛}，记作 $\mathscr{K}_X$（或 $\mathscr{K}_{X/k}$）。如 \S18.5.2 所述，若 $X$ 是射影的，则 $\mathscr{K}_X$ 就是\trSerreDuality{}的陈述中出现的那个\trCanonical{}层 $\omega_X$；这一点我们将在 \S29.4 中予以建立（见\mbox{\trDesideratum{}~29.1.1}）。

\textbf{21.5.B.（练习：$\mathscr{K}_X$ 的伴随公式）}\quad 设 $X$ 是光滑簇，$Z$ 是 $X$ 的一个光滑子簇。证明
\[
\mathscr{K}_Z\cong\mathscr{K}_X|_Z\otimes\det\mathscr{N}_{Z/X}.
\]
（提示：将\mbox{练习~14.2.P} 应用于\mbox{定理~21.2.31}。）特别地，由\mbox{练习~21.2.J}，若 $Z$ 的余维数为 $1$，则
\[
\mathscr{K}_Z\cong(\mathscr{K}_X\otimes\OO_X(Z))|_Z,
\]
这个式子更常用、也更紧凑的写法是 $\mathscr{K}_X(Z)|_Z$。伴随公式常被归纳地用于完全相交的情形；\mbox{练习~21.5.G} 即是一例。（伴随公式的其他版本见\mbox{练习~20.2.B(a)} 与 \S29.4。）

\textbf{21.5.4.}\ \emph{\trhDefinition。}\quad 若 $X$ 是射影（甚至是固有）的光滑 $k$-簇，则双有理不变量 $\mathrm{h}^0(X,\mathscr{K}_X)=\mathrm{h}^0(X,\Omega^n_{X/k})$ 具有特殊的重要性。它被称为\textbf{\trGeometricGenus}，记作 $\mathrm{p}_g(X)$。曲线的情形已在 \S21.3.1 中见过。若 $X$ 是一个未必光滑或射影的不可约簇，则「\trGeometricGenus」一词指的是与 $X$ 双有理等价的某个光滑射影簇 $X'$ 的 $\mathrm{h}^0(X',\mathscr{K}_{X'})$。（由\mbox{练习~21.5.A}，这与 $X'$ 的选取无关。）例如，若 $X$ 是 $k$ 上的一条既约不可约射影\trCurve，则\trGeometricGenus{}就是 $X$ 的\trNormalization{}的\trGenus。（但在更高的维数，这样的 $X'$ 是否存在并不清楚。一件非平凡的事实是：在\trCharacteristic{} $0$ 时这的确成立，而在正\trCharacteristic{}时这还是未知的；见\mbox{注记~22.4.6}。）

堪称奇迹的是：对一条复曲线而言，\trGeometricGenus{}与\trTopologicalGenus{}及\trArithmeticGenus{}相同。\mbox{练习~18.5.A} 证明了\trGeometricGenus{}与算术亏格相等（假定\mbox{奇迹~18.5.2}，即典范丛是\trSerre{}对偶化的）。我们马上就会在关于\trRiemann--\trHurwitz{}公式的讨论中把\trGeometricGenus{}与拓扑亏格联系起来（\mbox{练习~21.4.I}）。

\textbf{21.5.C.（练习）}\quad 设 $Z$ 是 $\mathbb{P}^3_{\bar{k}}$ 中一张次数为 $d$ 的正则曲面。计算 $Z$ 的\trGeometricGenus{} $\mathrm{p}_g(Z)$。证明 $\mathbb{P}^3_{\bar{k}}$ 中不存在有理的正则四次曲面（即与 $\mathbb{P}^2_{\bar{k}}$ 双有理等价，\mbox{定义~7.5.5}）。（这样的四次曲面是 K3 \emph{曲面}的一些例子；见\mbox{练习~21.5.I}。）

\textbf{21.5.5. 簇的几类重要类型：Fano、Calabi--Yau、一般型。}\quad 设 $X$ 是光滑射影 $k$-簇。若 $\mathscr{K}_X^\vee$ 是\trAmple{}的，则称 $X$ 是 \textbf{Fano} 的；若 $\mathscr{K}_X\cong\OO_X$，则称 $X$ 是 \textbf{Calabi--Yau} 的。（警示：Calabi--Yau 还有别的定义。）

\textbf{21.5.D.（练习）}\quad 对一切 $j\ge 0$，光滑射影 $k$-簇的\textbf{第 $j$ 个多亏格}（jth plurigenus）是 $\mathrm{h}^0(X,\mathscr{K}_X^{\otimes j})$。证明第 $j$ 个多亏格是双有理不变量。

（相比之下，当 $j<0$ 时 $\mathrm{h}^0(X,\mathscr{K}^{\otimes j})$ 不是双有理不变量。例如，你可以证明：对 $X=\mathbb{P}^2_k$ 与 $X=\mathrm{Bl}_p\mathbb{P}^2_k$ 而言，$\mathrm{h}^0(X,\mathscr{K}^{\vee})$ 的值有所不同。）

$X$ 的\textbf{\trKodaira{}维数}（Kodaira dimension），记作 $\kappa(X)$，追踪的是诸个多亏格的增长速率。它是使 $\mathrm{h}^0(X,\mathscr{K}_X^{\otimes j})/j^k$ 有界（当 $j$ 遍历正整数时）的最小的 $k$；唯一的例外是：若诸个多亏格 $\mathrm{h}^0(X,\mathscr{K}_X^{\otimes j})$ 全为 $0$，则我们说 $\kappa(X)=-1$。小平维数「总是存在」，并且它是诸第 $j$ 个「多典范有理映射」（pluricanonical rational maps）的像的维数中的最大者——这是一件非平凡的事实。后一事实蕴涵：小平维数是介于 $-1$ 与 $\dim X$ 之间（含两端）的整数。\mbox{练习~21.5.D} 表明小平维数是双有理不变量。若 $\kappa(X)=\dim X$，我们说 $X$ 是\textbf{一般型}（general type）的。（关于小平维数的更多信息，见 [Ii, \S10.5]。）

\textbf{21.5.E.（练习）}\quad 证明：若 $\mathscr{K}_X$ 是\trAmple{}的，则 $X$ 是一般型的。

\textbf{21.5.F.（练习）}\quad 证明：（域 $k$ 上的）光滑\trGeometricallyIrreducible{}射影\trCurve{}当其\trGenus{}为 $0$（相应地，$1$，大于 $1$）时是 Fano（相应地，Calabi--Yau，一般型）的。

\mbox{练习~21.5.F} 的「三分法」（trichotomy）在精神实质上正是如下现象的原因：曲线的行为依其所属的三个类而各不相同（见 \S19.6.1）。在某种意义上，这一三分法的一些重要侧面可以延伸到更高的维数，这也是我们要作这些定义的部分理由。现在我们在完全相交的情形探究这一三分法。

\textbf{21.5.G.（练习）}\quad 设 $X$ 是 $\mathbb{P}^N_k$ 中由次数为 $d_1,\dots,d_n$ 的超曲面给出的光滑完全相交，其中 $k$ 是代数闭的。
（a）证明 $\mathscr{K}_X\cong\OO(-N-1+d_1+\cdots+d_n)|_X$。
（b）证明：若 $-N-1+d_1+\cdots+d_n$ 为负（相应地，为零，为正），则 $X$ 是 Fano（相应地，Calabi--Yau，一般型）的。
（c）找出 $N$ 与 $d_1,\dots,d_n$ 的所有可能的取值，使得 $X$ 是维数至多为 $3$ 的 Calabi--Yau。留意这张清单有多小。

\textbf{21.5.H.（容易的练习）}\quad 光滑 Fano 曲面称为\textbf{\trDelPezzo{}曲面}（del Pezzo surface）（它是曲面的\trEnriques--\trKodaira{}分类的一部分）。证明下列曲面都是 del Pezzo 曲面：(a) $\mathbb{P}^2_k$；(b) $\mathbb{P}^1_k\times\mathbb{P}^1_k$；(c) $\mathbb{P}^3_k$ 中次数至多为 $3$ 的光滑曲面。（这张清单有一些冗余！del Pezzo 曲面的其余例子在\mbox{注记~27.3.7} 中有描述；其中的一个早在 \S19.8.3 中就已出现过。）

\textbf{21.5.I.（练习）}\quad 域 $k$ 上的一个 \textbf{K3 曲面}（K3 surface）是 $k$ 上一个固有、光滑、\trGeometricallyConnected{}的 Calabi--Yau 曲面 $X$，满足 $\mathrm{H}^1(X,\OO_X)=0$。（\trWeil{}曾写道：K3 曲面的命名是为了纪念\trKummer、\trKahler、\trKodaira{}三位数学家以及 K2 山峰，[BPV, p.~288]。）证明\mbox{练习~21.5.G(c)} 中维数为 $2$ 的（光滑）Calabi--Yau 完全相交都是 K3 曲面。

\textbf{21.5.6.}\ \emph{诱人的题外注记。}\quad 若你把切出完全相交 K3 曲面的那些超曲面的次数（见\mbox{练习~21.5.G(c)}），与切出完全相交典范曲线的超曲面的次数（见 \S19.8.2 与\mbox{练习~19.8.E}）作个比较，你会注意到一个引人注目的巧合。当然，这根本不是巧合。

\textbf{21.5.J.（练习：并非一切都能纳入这一三分法）}\quad 设 $C$ 是\trGenus{}大于 $1$ 的光滑射影不可约复曲线。证明 $C\times_\mathbb{C}\mathbb{P}^1_\mathbb{C}$ 既不是 Fano 的，也不是 Calabi--Yau 的，还不是一般型的。（提示：\mbox{练习~21.2.U}。）

\textbf{21.5.7.}\ \hard\ \textbf{\trKodaira{}消失。}\quad \trKodaira{}消失定理是代数几何中好几个领域的重要工具。我们出于文化修养的缘故陈述它，但不予证明（因此也不会用到它）。

\textbf{21.5.8.\ \trhTheorem（\trKodaira{}消失定理）。}\ ——\ \emph{设 $k$ 是\trCharacteristic{} $0$ 的域，而 $X$ 是光滑射影 $k$-簇。那么对任意\trAmple{}\trInvertibleSheaf{} $\mathscr{L}$，当 $i>0$ 时 $\mathrm{H}^i(X,\mathscr{K}_X\otimes\mathscr{L})=0$。}

\textbf{21.5.9.}\quad 对\trCharacteristic{}的限制是必要的：\trRaynaud{}给出了一个在正\trCharacteristic{}下小平消失失效的例子 [Ra]。原始证明的做法是：先证复解析的版本，再借\trSerre{}的 GAGA 定理 [Se3] 把它转移成一个复的代数陈述。一般的\trCharacteristic{} $0$ 情形可以归结为 $\mathbb{C}$ 的情形——任何这类归结常常被（有些含糊地）称作\trLefschetz{}原理（Lefschetz principle）的一次应用。（\trLefschetz{}原理的一个精确表述，连同适用于任意\trCharacteristic{}的短证明，见 [Ek]；一个推广见 [FR]；更多背景见 [MO90551]。）\trRaynaud{}后来给出了小平消失的一个戏剧性的代数证明，其中\emph{用到}了正\trCharacteristic（！）；见 [De1] 或 [Il]。

\textbf{21.5.K.（练习）}\quad 在 $\dim X=1$ 的情形以及对 $\mathbb{P}^n_k$ 证明小平消失定理。（这两种情形都不需要\trCharacteristic{} $0$ 的假设。）

\textbf{21.5.10.}\ \hard\ \textbf{初窥\trHodge{}理论。}\quad 不变量 $\mathrm{h}^j(X,\Omega^i_{X/k})$ 称为 \textbf{\trHodge{}数}（Hodge number）$\mathrm{h}^{i,j}(X)$。由\mbox{练习~21.5.A}，诸 $\mathrm{h}^{i,0}$ 是双有理不变量。我们很快就会看到（在\mbox{练习~21.5.O} 中），这一点并非对所有的 $\mathrm{h}^{i,j}$ 都成立。

\textbf{21.5.L.（练习）}\quad 设 $X$ 是光滑射影簇。假定\mbox{奇迹~18.5.2}（典范丛是\trSerre{}对偶化的），证明诸\trHodge{}数满足对称性 $\mathrm{h}^{p,q}=\mathrm{h}^{n-p,n-q}$。（\mbox{练习~14.2.O} 会派上用场。）

\textbf{21.5.M.（练习：射影空间的 \trHodge{}数）}\quad 证明：当 $0\le p=q\le n$ 时 $\mathrm{h}^{p,q}(\mathbb{P}^n_k)=1$，否则 $\mathrm{h}^{p,q}(\mathbb{P}^n_k)=0$。（提示：使用欧拉正合列（\mbox{定理~21.3.6}），并应用\mbox{练习~14.2.N}。）

\textbf{21.5.11.}\ \emph{\trhRemark{}：\trHodge{}菱形。}\quad 在 $k=\mathbb{C}$ 上，还会发生更多的奇迹。若 $X$ 是不可约光滑射影复簇，那么结果表明存在一个直和分解
\begin{equation*}\tag{21.5.11.1}\label{eq:hodge-decomposition}
\mathrm{H}^m(X,\mathbb{C})=\bigoplus_{i+j=m}\mathrm{H}^j(X,\Omega^i_{X/\mathbb{C}}),
\end{equation*}
从而 $\mathrm{h}^m(X,\mathbb{C})=\sum_{i+j=m}\mathrm{h}^{i,j}$，于是诸\trHodge{}数（纯代数的对象）给出了贝蒂数（Betti numbers；先验地看，这是拓扑信息）。此外，复共轭把 $\mathrm{H}^j(X,\Omega^i_{X/\mathbb{C}})$ 与 $\mathrm{H}^i(X,\Omega^j_{X/\mathbb{C}})$ 互换，由此得到
\begin{equation*}\tag{21.5.11.2}\label{eq:hodge-symmetry}
\mathrm{h}^{i,j}=\mathrm{h}^{j,i}.
\end{equation*}
（题外话：这条附加的对称性在\trCharacteristic{} $0$ 下普遍成立，但在正\trCharacteristic{}下可能失效；例如见 [Ig, p.~966]、[Mu1, \S III]、[Se4, Prop.~16]。）这是广阔而丰富的\trHodge{}理论这一主题的开端（想多知道一点，见 [GH1, \S0.6]；想多知道很多，见 [Vo]）。

若我们把诸\trHodge{}数写成一个菱形，其中 $\mathrm{h}^{i,j}$ 是第 $i+j$ 行内的第 $i$ 个条目，那么这个菱形便具有分别来自\trSerreDuality{}与复共轭的两个对称性。例如，一个不可约光滑射影复曲面的\trHodge{}菱形将具有如下形状：
\[
\begin{array}{ccccc}
 & & 1 & & \\
 & q & & q & \\
 \mathrm{p}_g & & \mathrm{h}^{1,1} & & \mathrm{p}_g \\
 & q & & q & \\
 & & 1 & &
\end{array}
\]
其中 $\mathrm{p}_g$ 是该曲面的\trGeometricGenus，而 $q=\mathrm{h}^{0,1}=\mathrm{h}^{1,0}=\mathrm{h}^{2,1}=\mathrm{h}^{1,2}$ 称为该曲面的\textbf{不规则性}（irregularity）。再举一例，由\mbox{练习~21.5.M}，$\mathbb{P}^n$ 的\trHodge{}菱形除竖直对称轴上的一列 $1$ 之外全为 $0$。

你不需要那两条未加证明的陈述 \eqref{eq:hodge-decomposition} 或 \eqref{eq:hodge-symmetry}，也能解出下面的习题。

\textbf{21.5.N.（练习）}\quad 假定\mbox{奇迹~18.5.2}（典范丛是\trSerre{}对偶化的），证明：域 $k$ 上一条光滑射影\trGeometricallyIrreducible{}、亏格为 $g$ 的曲线，其\trHodge{}菱形如下：
\[
\begin{array}{ccc}
 & 1 & \\
 g & & g \\
 & 1 &
\end{array}
\]

\textbf{21.5.O.（练习）}\quad 证明 $\mathbb{P}^1_k\times\mathbb{P}^1_k$ 的\trHodge{}菱形如下：
\[
\begin{array}{ccccc}
 & & 1 & & \\
 & 0 & & 0 & \\
 0 & & 2 & & 0 \\
 & 0 & & 0 & \\
 & & 1 & &
\end{array}
\]
把你得到的答案与 $\mathbb{P}^2_k$ 的\trHodge{}菱形（\mbox{练习~21.5.M}）作比较，证明 $\mathrm{h}^{1,1}$ 不是双有理不变量。

注意在上述两种情形中，$\mathrm{h}^{1,1}$ 都是\trPicard{}数 $\rho$（定义见 \S18.4.12）。一般地，$\rho\le\mathrm{h}^{1,1}$（\trCharacteristic{} $0$ 时见 [GH1, \S3.5, pp.~456--7]）。

\textbf{21.5.12.}\ \hard\ \textbf{题外：无穷小形变与自同构。}\quad 把这件事说精确超出了本书的范围，但若 $X$ 是一个簇，则 $\mathrm{H}^0(X,\mathscr{T}_X)$ 参数化 $X$ 的无穷小自同构，而 $\mathrm{H}^1(X,\mathscr{T}_X)$ 参数化 $X$ 的无穷小形变。作为一个例子，若 $X=\mathbb{P}^1$（在某个域上），则 $\mathscr{T}_{\mathbb{P}^1}\cong\OO(2)$（\S21.3.2），于是 $\mathrm{h}^0(\mathbb{P}^1,\mathscr{T}_{\mathbb{P}^1})=3$，而这恰好就是 $\mathbb{P}^1$ 的自同构群的维数（\mbox{练习~15.5.A}）。

\textbf{21.5.P.（练习）}\quad 用欧拉正合列（\mbox{定理~21.3.6}）计算 $\mathrm{h}^0(\mathbb{P}^n_k,\mathscr{T}_{\mathbb{P}^n_k})$，并将此与 $\mathbb{P}^n_k$ 的自同构群的维数（\mbox{练习~15.5.A}）作比较。

\textbf{21.5.Q.（练习）}\quad 证明 $\mathrm{H}^1(\mathbb{P}^n_A,\mathscr{T}_{\mathbb{P}^n_A})=0$。（于是射影空间不能形变，是「刚性」的。）

\textbf{21.5.R.（练习）}\quad 假定\mbox{奇迹~18.5.2}（典范丛是\trSerre{}对偶化的），对 $k$ 上一条亏格为 $g$ 的光滑射影\trGeometricallyIrreducible{}曲线 $C$，就 $i=0$ 与 $1$ 计算 $\mathrm{h}^i(C,\mathscr{T}_C)$。你应当注意到：亏格为 $0$、$1$ 与 $g>1$ 时，$\mathrm{h}^i(C,\mathscr{T}_C)$ 分别为 $0$、$1$ 与 $3g-3$；做完之后，重读 \S19.8.4。
